\section{Summary of findings}

This thesis presented \lecseg{}: a benchmark, evaluation framework, and
diagnostic pipeline for low-resource lecture-video topic segmentation.
Two findings stand out from the full experimental programme.

The first concerns \emph{where the error actually lives}: an oracle experiment
showed that near-perfect candidate boundaries already exist in the generated
pool at 96.8\% recall ($P_k=0.0172$). The $\Delta P_k \approx 0.34$ gap between
that oracle and the best deployable system is, within the framework we evaluated,
predominantly a ranking problem rather than a detection problem. Generating
more or better candidates is not the bottleneck; choosing the right ones is.

The second finding concerns \emph{why multimodal signals keep failing}: prosody,
discourse markers, BERTopic, and GPT-2 perplexity all over-segment relative to
YouTube chapter granularity, not because those signals are uninformative but
because they detect transitions at a finer discourse granularity than the
editorial decisions creators make when writing chapter titles. CLIP visual
embeddings are the exception --- slide scene changes are themselves editorial
artifacts at approximately chapter granularity, which explains why CLIP succeeds
where the acoustic signals do not.

The quantitative results are in Table~\ref{tab:lecseg_result_variants} and
the statistical tests in Table~\ref{tab:significance_full}; the research
questions are answered directly below.

\section{Research questions answered}
\label{sec:rqs_answered}

\textbf{RQ1} (most reliable text-embedding and cross-model strategies):
Conservative cross-model scoring using BGE and E5 embeddings is the most
reliable single strategy, achieving $P_k=0.3713$ and WD$=0.3764$ with
statistical support over all classical, neural, and multimodal baselines.
BGE-divisive segmentation is the strongest single-model text baseline
at $P_k=0.3884$.

\textbf{RQ2} (per-video selector improvement beyond fixed strategy):
Yes, with qualifications.
The leave-one-video-out balanced selector improves mean $P_k$ to 0.3588,
statistically significant over the BGE-divisive baseline.
In practical terms, a $P_k$ of 0.3588 means the system generally identifies
the correct chapter region but remains conservative in recovering exact creator
boundary locations, typically placing boundaries within the right neighbourhood
rather than at the precise transition sentence.
The selector's $P_k$/WD advantage over the cross-model method does not reach
significance, and it degrades substantially ($P_k=0.4012$) when entire domains
are excluded from training.
The selector is therefore a useful low-resource operating point, not a
robust replacement for the fixed cross-model method.

\textbf{RQ3} (hierarchical annotation and nesting constraints):
Yes, the two-level annotation format with enforced nesting provides a richer
inspection surface than flat chapter labels.
The nesting enforcement produces valid hierarchical outputs by construction.
However, the hierarchical segmenter achieves higher $P_k$ than the
chapter-only methods (0.4118 vs.\ 0.3884), indicating that solving the
two-level task simultaneously remains harder than the single-level variant.

\textbf{RQ4} (which auxiliary signals help or hurt):
CLIP visual embeddings help: CLIP alone reaches $P_k=0.3958$ and CLIP plus
text fusion reaches $P_k=0.3740$, approaching the cross-model text result.
All other auxiliary signals tested, including prosody, discourse markers,
BERTopic, GPT-2 perplexity, and local LLM refinement, produce worse $P_k$/WD
than the BGE-divisive baseline when used as primary or fused signals.
The consistent explanation is granularity mismatch: these signals detect
sub-chapter transitions rather than editorial chapter boundaries.

\textbf{RQ5} (statistically significant improvements under the LECSEG-30 protocol):
Two improvements are statistically significant after Holm correction:
the cross-model conservative method over BGE-divisive ($\Delta P_k=-0.0171$,
$p<0.05$), and the balanced selector over BGE-divisive ($\Delta P_k=-0.0296$,
$p<0.05$).
The selector's $P_k$/WD margin over the cross-model method is not significant,
but its Boundary Similarity and F1@2 improvements over that method are.

\section{Contributions to the Field}

The contributions below are ordered by scientific significance:

\begin{enumerate}
  \item \textbf{\dataset{} benchmark.}
    30 videos, 32.52 hours, 419 creator chapter boundaries, and 904 reviewed
    subtopic labels.
    This is the primary contribution: it makes the low-resource lecture
    segmentation problem measurable and reproducible under a shared protocol.

  \item \textbf{Candidate-selection bottleneck finding.}
    Oracle analysis places the dominant source of error in selection, not
    generation (candidate oracle $P_k=0.0172$, 96.8\% recall; deployable best
    $P_k=0.3588$). This relocates future research toward boundary ranking.
    See Section~\ref{sec:oracle_gap_discussion} for the full analysis.

  \item \textbf{Granularity mismatch finding.}
    Non-embedding signals (prosody, BERTopic, perplexity, discourse markers)
    all hurt $P_k$/WD because they detect sub-chapter discourse transitions.
    CLIP is the principled exception: slide-scene changes share the editorial
    granularity of creator chapters. This finding transfers directly to any
    future system targeting editorial-chapter-level segmentation.
    See Section~\ref{sec:ablations} for ablation details.

  \item \textbf{Hierarchical annotation framework.}
    A two-level annotation protocol with independent human double annotation
    and a nesting-enforcement mechanism.
    Inter-annotator agreement ($\kappa_{\text{chapter}}=0.5351$,
    $\kappa_{\text{subtopic}}=0.4257$) confirms moderate human agreement
    on the subtopic task; the chapter figure is inflated because both
    annotators observe the same creator-provided YouTube metadata.

  \item \textbf{Reproducibility package.}
    Open pipeline covering transcription, segmentation (four integrated
    components and seven baselines), evaluation, and statistical testing.
\end{enumerate}

\subsection*{Scientific knowledge contributions}

The five contributions above produce four transferable knowledge claims:
(i)~boundary selection is the primary bottleneck in low-resource lecture
segmentation;
(ii)~granularity mismatch is the systematic explanation for multimodal signal
failures, not signal uninformativeness;
(iii)~CLIP is the only tested non-text signal that operates at editorial-chapter
scale, and for a principled reason;
(iv)~the negative results (OCR, prosody, discourse markers, BERTopic, perplexity,
and LLM refinement all failing to improve chapter-level $P_k$/WD) have direct
replication value for future projects.

\subsection*{Scope of the claims}

Dense text embeddings outperform classical and neural segmenters; candidate
selection, not generation, is the dominant error source; prosodic and discourse
signals over-segment relative to creator chapter granularity; CLIP is the one
non-text signal that aligns with editorial chapter scale.
That is the scope. What \dataset{} does \emph{not} show: that the segmentations
are pedagogically optimal, that students benefit measurably, or that findings
transfer to annotation standards other than creator chapters.

\section{Limitations}

\begin{itemize}
  \item \textbf{Corpus size and balance.}
    \dataset{} is a seed benchmark of 30 videos. The manifest covers five
    domains, but the distribution is not perfectly balanced.

  \item \textbf{Silver chapter labels.}
    Creator-provided YouTube chapters are useful but imperfect. They may reflect
    production choices rather than strictly pedagogical topic boundaries.

  \item \textbf{Conservative operating point.}
    The best $P_k$/WD method deliberately under-predicts boundaries to minimise
    false positives. This results in low diagnostic F1, which is expected and
    not a flaw: approximate segment structure and exact boundary recovery are
    different objectives, and the system is optimised for the former.
    A deployment that requires precise timestamp recovery would need a different
    operating point, likely at the cost of higher $P_k$.

  \item \textbf{Multimodal signal reliability.}
    Acoustic (pause/pitch), language-model perplexity, and BERTopic methods
    all produce worse $P_k$ than BGE-divisive, reflecting a granularity
    mismatch. The exception is CLIP visual embeddings: CLIP alone reaches
    $P_k=0.3958$, and CLIP plus text fusion reaches $P_k=0.3740$, below the
    cross-model text method ($P_k=0.3715$) but better than BGE-divisive
    ($P_k=0.3884$). Visual scene changes in lecture slides carry chapter-level
    granularity and are a promising signal for future work.

  \item \textbf{LLM refinement.}
    The local LLM component is implemented, but the current verified metrics do
    not justify claiming large boundary-placement gains from LLM refinement.
\end{itemize}

\section{Threats to validity}
\label{sec:threats_validity}

The following table summarises the principal threats and their estimated
severity.

\begin{table}[h]
  \centering
  \caption{Threat-to-validity overview for \lecseg{}.}
  \label{tab:threats}
  \small
  \begin{tabular}{p{4.5cm}lp{5.5cm}}
    \toprule
    Threat & Severity & Primary impact \\
    \midrule
    Creator chapters as reference & High & Construct validity:
      benchmark measures creator metadata, not pedagogical truth \\
    Small and imbalanced corpus & High & External validity:
      30 videos limit generalisation claims \\
    Single-domain annotator pool & Low--Medium & Annotation validity:
      both annotators share cultural and academic background; agreement
      may not generalise to other annotator groups \\
    Gold boundary count in evaluation & Medium & Internal validity:
      deployment realism; real systems cannot know $k$ in advance \\
    Post-hoc method search & Medium & Internal validity:
      multiple testing risk mitigated by Holm correction \\
    Domain imbalance (4 Math videos) & Medium & Selector stability on
      small domain splits \\
    ASR quality not formally measured & Low--Medium & Upstream noise
      propagation from Whisper output \\
    \bottomrule
  \end{tabular}
\end{table}

\paragraph{Construct validity.}
The benchmark evaluates a system's ability to reproduce creator-provided
navigation metadata, not pedagogically verified topic boundaries.
As discussed in Section~\ref{sec:youtube_validity}, creators may place chapter
markers for editorial or discoverability reasons that do not coincide with
genuine topic shifts.
A system achieving low $P_k$ is correctly predicting creator chapter structure;
whether this predicts learning benefit requires an independent user study.

\paragraph{External validity.}
\dataset{} has 30 videos across five domains. Results should be treated as
findings on this specific benchmark rather than as universal laws of lecture
segmentation. The leave-domain-out selector degradation ($P_k=0.4012$) is
direct evidence that performance is sensitive to domain-matched training data;
with more videos the picture may shift materially.
A ten-video external spot-check across five domains
(Section~\ref{sec:external_validation}) found BGE-divisive
consistent (external mean $P_k=0.389$ vs $0.388$ on \dataset{},
$\Delta=0.001$), while the cross-model method degraded somewhat
($P_k=0.421$ vs $0.371$, $\Delta=0.050$).
Confounds include yt-dlp auto-captions vs.\ Whisper transcripts
and an external pool skewed toward 3Blue1Brown explainer videos.

\paragraph{Annotation validity.}
Subtopic labels were produced by human annotators and verified through
independent double annotation. The inter-annotator agreement figure
($\kappa_{\text{subtopic}}=0.4257$) represents genuine human agreement
on a subjective discourse-segmentation task, which is in the moderate range
and appropriate for this type of annotation.
The chapter-level $\kappa=0.5351$ is inflated because both annotators derive
chapter boundaries from the same YouTube creator metadata; the subtopic figure
is the more meaningful validity measure.

\paragraph{Internal validity.}
Methods in the controlled evaluation receive the gold segment count.
This isolates boundary placement, which is standard practice in topic
segmentation, but it is not deployment-realistic.
Boundary count error is reported as a diagnostic
(Section~\ref{sec:modern_metrics}), and automatic count prediction is listed
as future work.
The use of many candidate methods and operating points is mitigated by
Holm correction over seven pre-specified comparisons, but the selector
result should still be read as a low-resource operating point rather than
a universal deployment guarantee.

\paragraph{Metric validity.}
$P_k$ and WindowDiff correlate with navigation usefulness under the assumption
that students find approximate-region boundary placement acceptable.
This assumption is plausible but untested; a user study measuring actual
navigation speed and retention would be needed to verify it empirically.

\section{Closing the oracle gap}
\label{sec:oracle_gap_discussion}

The oracle analysis is the clearest quantitative finding this thesis produces.
With perfect candidate selection from the generated pool, $P_k=0.0172$ is
reachable at 96.8\% recall.
The gap between this oracle and the best deployable system
($\Delta P_k \approx 0.34$) indicates that, within the evaluated framework,
the candidate pool already contains the information needed for near-perfect
segmentation; the bottleneck appears to lie in correctly identifying which
candidates to accept rather than in generating more candidates.
This framing is the most actionable scientific output of the thesis: it shifts
the research problem from \emph{how to generate better candidates} to
\emph{how to rank candidates more reliably}.

Table~\ref{tab:oracle_gap} summarises the relevant operating points.
TreeSeg (Gklezakos et al.\ 2024) reports $P_k=0.367$ on TinyRec (21
self-recorded lectures), and \lecseg{} achieves $P_k=0.3588$ with the balanced
selector on \dataset{}.
These figures are indicative rather than directly comparable because the
benchmarks, annotation sources, and video types all differ.

Three concrete reasons favour TreeSeg in any hypothetical direct comparison.
First, it uses overlapping utterance-block embeddings, which provide
richer contextual evidence per position than the single-sentence windows used
in \lecseg{}.
Second, TinyRec uses manually verified annotations rather than creator
chapter metadata, potentially making the reference labels more consistent.
Third, \lecseg{}'s selector incurs leave-domain-out degradation that would
not surface under same-domain evaluation.
None of these diminish the \lecseg{} findings; they define the comparison boundary.

\begin{table}[h]
  \centering
  \caption{Oracle gap analysis on LECSEG-30.}
  \label{tab:oracle_gap}
  \begin{tabular}{lrr}
    \toprule
    Operating point & $P_k\downarrow$ & WD$\downarrow$ \\
    \midrule
    BGE-divisive baseline & 0.3884 & 0.3956 \\
    Cross-model conservative (best global) & 0.3713 & 0.3764 \\
    Balanced selector (best deployable) & 0.3588 & 0.3739 \\
    TreeSeg paper (TinyRec, indicative only) & 0.367 & --- \\
    Per-video oracle (not deployable) & 0.2980 & 0.3280 \\
    Candidate oracle $\tau=2$ (upper bound) & 0.0172 & 0.0198 \\
    \bottomrule
  \end{tabular}
\end{table}

The concrete research directions that could close this gap are:

\begin{enumerate}
  \item \textbf{Supervised boundary-level ranker.}
    The video-level method selector shows that per-video method choice
    matters, but it does not rank individual boundary candidates.
    A ranker trained on candidate-level features (local cosine contrast,
    cross-model agreement, pause duration, slide-change signal, and OCR
    evidence) would directly address the selection bottleneck.

  \item \textbf{Expanded annotation.}
    The current selector is trained leave-one-out on 30 videos.
    Expanding to 50--100 annotated lectures would provide 2--3$\times$ more
    per-fold evidence and may substantially stabilise domain-specific selector
    behaviour.

  \item \textbf{Shared-dataset evaluation.}
    Running TreeSeg's original code on \dataset{} (or publishing \dataset{} as
    a community benchmark) would convert the current indicative comparison
    into a reproducible head-to-head result.

  \item \textbf{CLIP-centred multimodal fusion.}
    CLIP is the only non-text signal that improves $P_k$/WD.
    A tighter integration of CLIP slide-scene evidence within the cross-model
    scoring framework, rather than as a secondary fusion component, could
    provide a more reliable visual prior for chapter-level transitions.
\end{enumerate}

\section{Lessons learned}
\label{sec:lessons}

Several lessons from this project are likely to carry over to related work.

\paragraph{Segmentation lessons.}
\begin{enumerate}
  \item \textbf{Dense text embeddings dominate at creator chapter granularity.}
    BGE and E5 cross-model scoring outperforms every multimodal combination
    tested.
    Semantic paragraph-level structure is the dominant signal for
    editorial-level chapter boundaries; fine-grained acoustic and lexical
    transitions are counterproductive.

  \item \textbf{Auxiliary signals must be calibrated to the reference
    granularity, not just to transition detection.}
    Prosody and discourse markers detect real transitions, but those transitions
    are an order of magnitude finer than creator chapters.
    ``Granularity mismatch'' is a more precise diagnosis than
    ``signal uninformative.''

  \item \textbf{The generation/selection distinction is the key design choice.}
    Near-perfect boundaries are present in the candidate pool at 96.8\% recall.
    Future systems should invest in selection and ranking, not in generating
    additional candidate positions.

  \item \textbf{CLIP is the only tested non-text signal that aligns with
    editorial chapter granularity.}
    This makes it the most valuable non-text signal to integrate next, and the
    principled reason is the shared editorial origin of slide transitions and
    creator chapters, not a coincidence in the data.
\end{enumerate}

\paragraph{Benchmark and annotation lessons.}
\begin{enumerate}
  \setcounter{enumi}{4}
  \item \textbf{Small benchmarks enable failure analysis that large benchmarks
    cannot.}
    A 30-video, fully auditable corpus makes it possible to inspect every
    failure case, trace every metric value to its source video, and report
    negative results honestly.
    Large-scale leaderboards measure performance; small auditable benchmarks
    measure understanding.

  \item \textbf{Creator chapters and pedagogical boundaries are different
    constructs.}
    Using creator chapters enables scale and reproducibility but couples the
    benchmark to editorial decisions.
    Any system trained or evaluated against creator chapters should be read
    as modelling navigation metadata rather than educational structure.

  \item \textbf{Negative results have replication value.}
    Documenting that prosody, discourse markers, BERTopic, OCR, and local LLM
    refinement all fail to improve chapter-level $P_k$/WD prevents future
    projects from repeating these experiments without motivation.
    Most published theses omit failures; their inclusion here is a deliberate
    scientific choice.
\end{enumerate}

\section{Closing remarks}

Lecture segmentation is far from solved, and this thesis does not pretend
otherwise.
What it contributes is a reproducible point of measurement and a clear
diagnosis of where the remaining difficulty lies.

The two central findings --- that the dominant error is in boundary selection
not detection, and that many signals fail because they operate at the wrong
discourse granularity --- together define the research agenda for any
successor system.
A future design that addresses both, with a learned boundary ranker and
signals explicitly calibrated to editorial chapter scale, has a clear,
measurable target: close the $\Delta P_k \approx 0.34$ oracle gap.

There is also a broader point about benchmark scale worth making explicitly.
Large benchmarks optimise for performance measurement: they reveal which
system produces the best aggregate number on a large diverse test set.
Small, auditable benchmarks optimise for \emph{understanding}: they reveal
why systems fail, which signals over- or under-segment, and where exactly
the performance ceiling lies.
\dataset{} was designed for the second purpose.
The granularity mismatch finding and the oracle-gap diagnosis are both results
that require the ability to inspect every prediction on every video ---
something a 200,000-video benchmark makes impractical.
The value of \lecseg{} is not in competing with large-scale systems
but in providing the kind of controlled, transparent analysis those systems
cannot easily afford.

The \dataset{} corpus, annotation labels, evaluation code, and result files
are open, so any future method can be evaluated against the same 30 videos
under the same protocol.
That reproducibility is what allows the findings here to be built upon rather
than merely noted.

If there is one sentence in this thesis that carries the most weight, it is:
\emph{candidate generation is largely solved; candidate selection is the
bottleneck.}
That observation, supported by the oracle-gap analysis, tells the next
researcher exactly where to invest effort --- and it could only have been
discovered by building the benchmark first.

\section{Recommendations for Future Work}

\subsection{Robust candidate and method selection}

The strongest immediate extension is a sequence-aware candidate selector.
Current oracle and method-portfolio experiments show that plausible boundaries
are often present in the candidate pool, but the system does not reliably pick
the right subset for each video. A better selector should combine local
semantic contrast, cross-model agreement, pause/shot/OCR evidence, position
priors, and global duration constraints while optimising directly against
$P_k$, WindowDiff, and tolerance-F1.

\subsection{Independent validation of chapter references}

The most important dataset extension is not simply adding more videos; it is
checking whether creator-provided YouTube chapters agree with independent
human judgements. A focused validation study should select 8--10 lectures and
ask two annotators to mark chapter-level boundaries without seeing the YouTube
timestamps. The study should report human--human agreement, human--YouTube
agreement, and system--human agreement under F1@2, F1@5, and time-tolerant
metrics. This would quantify how reliable the current creator chapters actually
are as reference boundaries.

\subsection{Expanded annotator pool and domain diversity}

The current inter-annotator study involves annotators from a shared academic
background. Future work should recruit annotators from different disciplines
and experience levels to test whether subtopic agreement generalises beyond
the current pool. Annotating 8--10 additional videos with three or more
independent annotators would give a Krippendorff's $\alpha$ estimate alongside
the existing Cohen's $\kappa$, and would allow pairwise agreement patterns to
be analysed rather than collapsed to a single mean figure.

\subsection{External and Math-focused validation}

The current selector fails most clearly on Mathematics, where the dataset has
only four videos. A targeted extension should add 8--10 Math lectures with
creator chapters and re-evaluate the final methods without retuning. If Math
performance improves, the likely cause was training scarcity; if it stays weak,
the problem probably runs deeper, involving mathematical notation, stable
vocabulary, or gradual conceptual transitions. A separate 5--10 video external
test set from new channels should also be used to measure generalisation beyond
\dataset{}.

\subsection{Reliability-aware prosody and visual integration}

Prosodic features (pause duration, pitch reset) and visual streams (shot
boundaries, slide OCR) are available in the pipeline, but current ablations
show that acoustic and linguistic sub-sentence signals can hurt
$P_k$/WindowDiff when their granularity does not match chapter-level editorial
decisions. The exception is CLIP visual embeddings: CLIP alone achieves
$P_k=0.3958$ and CLIP plus text fusion reaches $P_k=0.3740$, approaching the
cross-model text method ($P_k=0.3715$). Visual scene changes in lecture slides
appear to carry chapter-level granularity. Future work should prioritise CLIP
visual integration within the cross-model scoring framework, and apply
quality-gating only to the noisier acoustic and OCR modalities.

\subsection{Streaming and online inference}

The current pipeline processes a complete video offline.
A streaming variant would segment lecture content as it is being recorded,
enabling real-time chaptering for live classes.
This requires adapting the sliding-window fusion to operate on a causal context
window and replacing the global threshold computation with an exponential moving
average of past gap scores.
The target latency budget for live chaptering is under 30~seconds of delay.

\subsection{Cross-lingual extension}

\lecseg{} currently targets English lectures.
Extending to other languages requires: (1)~replacing \texttt{all-mpnet-base-v2}
with a multilingual text encoder such as LaBSE or \texttt{multilingual-e5-large};
and (2)~using Whisper's multilingual mode for transcription.
Suitable evaluation data could come from non-English MOOCs (e.g.,
Coursera's French and Spanish tracks) or from university recordings in South
and South-East Asian languages.

\subsection{End-to-end fine-tuning}

The current architecture uses frozen pre-trained encoders.
Fine-tuning the text encoder jointly with the boundary predictor on a larger
annotated corpus (e.g., constructed by scaling the \dataset{} annotation
procedure to 200+ videos) could yield further gains, at the cost of reduced
generalisability.
A natural training objective is a soft version of the tolerance-$F_1$ metric,
differentiable via a sigmoid boundary-probability output.

\subsection{Larger and domain-tuned local LLMs}

Replacing Llama-3.1-8B with a 70B parameter model or a domain-specific
fine-tune (e.g., fine-tuned on lecture transcripts) could improve title
quality and boundary-shift accuracy, particularly for highly technical content.
Quantised models (GGUF 4-bit) make 70B inference feasible on a single 24\,GB
GPU with acceptable latency (around 5~seconds per boundary refinement call).

\subsection{Multi-speaker lectures and panels}

Panel discussions and multi-instructor courses introduce speaker turns as an
additional boundary signal.
Speaker diarisation (e.g., using pyannote.audio) could provide a speaker-change
feature that complements text embeddings, particularly for sections where the
topic does not change but the speaker does.

\subsection{User study on learning outcomes}

All metrics in this thesis are information-retrieval proxies for segmentation
quality.
A controlled user study would directly measure whether learners locate
target concepts faster, and whether their retention improves, when using
\lecseg{}-generated chapters compared to flat video scrubbing or manually
authored chapters.
Such a study could also reveal which granularity (chapter vs.\ subtopic)
is most useful for different learning tasks (review vs.\ first-watch).

\subsection{Public benchmark and model release}

A complete public release should include the manifest, video identifiers,
creator chapter references, reviewed subtopic labels, preprocessing scripts,
evaluation scripts, environment lock files, result-regeneration commands, and
dataset/model cards. The current repository already contains the major
components; the remaining work is release hygiene, privacy review, and a
single documented reproduction command for the final thesis tables.

